\subsection{The central kernel over the universal quotient is finite}

The following calculation uses the ordinary discrete groups
first, so no rational continuous-cohomology identification
for a profinite group is being assumed.  Put
$B_0=C_2*C_3$, let $F_0$ be the specified index-six free
subgroup, and take the six discrete maps and kernels
corresponding to (B.R1).  Denote their intersection by
$K_0$, and put $U_0=B_0/K_0$, $J_0=F_0/K_0$.
The discrete version of the transfer and commutator proof
of B.2.10 gives the same finite-index subgroup
$J_0\subseteq B_0^6$.  Let $L$ be the discrete kernel of
the marked quotient $B_0\to\mathcal T$ of (B.R7).
Then $L^6\subseteq J_0$, and
\[
 U_0/L^6=E,\qquad |E|=11943936=2^{14}3^6.                       \tag{B.S4}
\]
The subgroup $L$ is torsion free: the two finite factors
of $B_0$ retain their orders in $\mathcal T$, and every
torsion element of their free product is conjugate into
a factor.  It is a free group of rank $25$, by its index
$144$ and the Euler characteristic $\chi(B_0)=-1/6$.

\begin{result}{B.2.15}{Proved}{Finiteness of the actual universal central kernel}
For the universal quotient $U$ of B.2.13, the actual group
\[
 M=\ker(U^{\mathrm c}\longrightarrow U),\qquad
 U^{\mathrm c}=B/\overline{[B,K]},                              \tag{B.S5}
\]
is finite and is annihilated by $2^{14}3^6$.  In particular
its prime divisors are among two and three.
\end{result}

\begin{proof}
We first prove the exact rational homology vanishing
\[
                         H_2(U_0,\mathbb Q)=0.                 \tag{B.S6}
\]
For a finite-index normal subgroup, transfer and averaging
identify rational homology of the ambient group with the
coinvariants of the subgroup homology.  Applied to
$L^6\subseteq J_0\subseteq U_0$, this gives
\[
\begin{aligned}
 H_2(J_0,\mathbb Q)
   &\cong\left(
        \bigoplus_{i<j}V_i\otimes V_j\right)_{J_0/L^6},
       & V_i&=H_1(L,\mathbb Q),\\
 H_2(U_0,\mathbb Q)&\cong H_2(J_0,\mathbb Q)_{S_3}.
\end{aligned}                                                   \tag{B.S7}
\]
The first displayed sum is the Kunneth formula for a
product of six free groups; each individual factor has
zero second homology.  These averaging statements also
follow from the homology spectral sequence of the finite
quotient: all its positive-degree rational homology
vanishes, since division by its order averages the bar
complex.

We determine the pair coinvariants in (B.S7) from their
actual common quotients.  For two epimorphisms from $F_0$,
their joint image is the fibre product over the quotient
by the product of their kernels.  The table (B.R5)
therefore shows that, from coordinate $1$, the common
quotient with coordinate $2$ is $S_3=B_0/T_2$, with
coordinate $5$ it is $A_4=B_0/T_3$, and with each of the
other three coordinates it is trivial.  The remaining
pairs are obtained by the actual source $S_3$ permutation
of the six coordinates.  This yields two exceptional
three-edge matchings; the other nine pairs have trivial
common quotient.

For a pair with trivial common quotient, the joint action
on the two $V$ factors is the full independent
$B_0\times B_0$ action.  Its coinvariants vanish, because
$V_{B_0}=H_1(B_0,\mathbb Q)=0$.  For an exceptional
pair with common quotient $C=S_3$ or $A_4$, first taking
coinvariants by the independent kernels gives
$H_1(T_2,\mathbb Q)$ or $H_1(T_3,\mathbb Q)$ in each
factor.  This follows again by finite-index averaging,
now for $L\subseteq T_2$ or $L\subseteq T_3$.

These two representations have an explicit peripheral
description.  The kernels $T_2,T_3$ are respectively the
three-cusped and four-cusped sphere groups.  Indeed their
indices in the modular orbifold group are six and twelve,
the cusp image has order two and three respectively,
and the Euler characteristics are $-1$ and $-2$.
There are consequently three and four cusps and genus
zero in each case.  Positive cusp classes identify their
homology with the augmentation representations
$W_3=\mathbb Q^3/\langle(1,1,1)\rangle$ of $S_3$ and
$W_4=\mathbb Q^4/\langle(1,1,1,1)\rangle$ of $A_4$.
Both permutation actions are doubly transitive.  Their
tensor-square invariant spaces are one-dimensional,
spanned by the symmetric tensor $\Omega$ obtained from
the usual invariant dot product.  Equivalently, the two
orbits on ordered pairs give two invariants in the full
permutation tensor square, and removing the constant
summand leaves exactly one.

We retain the precise common-quotient twists and the
Koszul sign in the action of a matching stabilizer.
For the pair $(1,2)$, the common-quotient automorphism
$\theta$ interchanges $s$ and $\tau$ in $S_3$.  It is
conjugation by the transposition $s\tau s$; its action
$P$ on $W_3$ is orthogonal and satisfies $P^2=1$.
The invariant tensor for the twisted diagonal action is
$\Omega_P=(1\otimes P)\Omega$.  The source element
$s$ interchanges the two actual coordinates without an
additional inner map.  On second homology it therefore
acts by $-\operatorname{Swap}$.  Since $P^t=P$, this
takes $\Omega_P$ to $-\Omega_P$.

For the pair $(1,5)$, the actual common-quotient
automorphism of $A_4$ is
$\theta(s)=s$, $\theta(\tau)=s\tau^{-1}$.
Thus $\theta^2=\operatorname{Ad}_s$ on the displayed
generators.  Its permutation $P$ of the four cusps is
orthogonal and satisfies $P^2=\rho(s)$ on $W_4$:
the equality follows from $\theta^2=\operatorname{Ad}_s$
and the faithful four-point action, whose permutation
centralizer is trivial.  Here the matching stabilizer
is the source $\tau=sr$.  On the two actual coordinates
its action is
\[
                 (v,w)\longmapsto(w,svs^{-1}),
\]
because $\tau^2=x^{-1}$ and $\varphi_1(x^{-1})=s$.
Its second-homology action is consequently
$-(1\otimes\rho(s))\operatorname{Swap}$.  Applied to
$\Omega_P$, it is
$-(1\otimes\rho(s)P^t)\Omega=-\Omega_P$, since
$\rho(s)=P^2$ and $P^t=P^{-1}$.

Thus the stabilizer of each exceptional matching acts
by minus one on its only surviving rational line.
Taking the final $S_3$ coinvariants in (B.S7) kills both
lines, proving (B.S6).

The discrete group $U_0$ is finitely presented: it is a
finite extension of $J_0$, which has finite index in a
finite product of virtually free finitely presented
groups.  Its integral second homology is therefore
finitely generated.  Equation (B.S6) makes it finite.
Moreover transfer from $U_0$ to its subgroup $L^6$ has
target the torsion-free group $H_2(L^6,\mathbb Z)$ and
therefore vanishes on this finite group.  Composing
transfer with inclusion is multiplication by the index
$|E|$.  Hence $|E|$ annihilates $H_2(U_0,\mathbb Z)$.

The source abelianization $B_0^{\mathrm{ab}}\to
U_0^{\mathrm{ab}}$ is an isomorphism $C_6$, by the
same two-coordinate argument preceding (B.S1).
Also $H_2(B_0,\mathbb Z)=0$: the classifying space of
the free product is the wedge of the two cyclic
classifying spaces, whose second integral homology
vanishes.  The low-degree homology sequence for
$K_0\to B_0\to U_0$ consequently gives
\[
        K_0/[B_0,K_0]\cong H_2(U_0,\mathbb Z).                 \tag{B.S8}
\]
In particular this actual discrete central kernel is
finite and annihilated by $|E|$.

Finally the completion comparison retains the same
kernel.  The completion of $J_0$ is the open subgroup
$J\subseteq B^6$, because $J_0$ has finite index in
$B_0^6$ and carries its full profinite topology.  Both
$\widehat U_0$ and $U$ are extensions of this $J$ by
the same quotient $S_3$; the map induced from $B_0$
is an isomorphism on this subgroup and quotient, hence
an isomorphism.  Thus $K$ is the closure of $K_0$ in
$B$.  Continuity and density give
$\overline{[B,K]}=\overline{[B_0,K_0]}$.  Consequently
$U^{\mathrm c}$ is the profinite completion of
$B_0/[B_0,K_0]$, and its kernel $M$ is an image of the
finite group in (B.S8).  This proves its finiteness
and the stated annihilator without assuming that
every element of the discrete central kernel survives
completion.
\end{proof}
